\section{5. Interfaces and data contracts}

% Long literal names in this section: let lines stretch instead of sticking out.
\sloppy

Each contract says what is sent, in which units, how often, when it counts as
stale, and what happens when the other side does not answer. The same
contracts, with every field, are kept next to the code in
\code{docs/interfaces.md}.

Conventions for every interface:
\begin{itemize}
\tightlist
\item \textbf{Rooms} are keyed \code{<level>/<room>}, for example
  \code{level0/A109}. Room names repeat between floors.
\item \textbf{Model time} is RFC~3339 as occupancysim's clock gives it, for
  example \code{2026-09-21T10:15:00Z}. Real time is used only for network
  timeouts.
\item \textbf{Run id.} Every process takes \code{RUN\_ID} from its environment
  and puts it in every record it publishes.
\item Values written to BuildSim are \textbf{strings} (BuildSim's rule). Values
  on MQTT and on our own HTTP APIs are JSON numbers.
\item MQTT is QoS~1, not retained except the model time, with a login per
  service. The broker's access list decides who may read and write each topic.
\end{itemize}

\begin{longtable}[]{@{}L{0.06\linewidth}L{0.30\linewidth}L{0.19\linewidth}L{0.24\linewidth}L{0.09\linewidth}@{}}
\toprule
\textbf{ID} & \textbf{Interface} & \textbf{Writer} & \textbf{Readers} & \textbf{Protocol} \\
\midrule
\endhead
IF-1 & Model clock, people per minute & occupancysim; physics relays the time & physics; actuators & REST, MQTT \\
IF-2 & Truth API & physics & sensor devices only & REST \\
IF-3 & Sensor readings in BuildSim & each sensor device & viewer; controller and damper as fallback & REST \\
IF-4 & Observations \code{obs/\ldots} & each sensor device & controller, dampers, storage, viz & MQTT \\
IF-5 & Truth \code{truth/\ldots} & physics & storage only & MQTT \\
IF-6 & Decisions \code{decision/\ldots} & controller & storage, viz & MQTT \\
IF-7 & Actuator events \code{event/\ldots} & actuator devices & storage, viz & MQTT \\
IF-8 & Actuator command & controller & one actuator device & REST \\
IF-9 & Actuator state in BuildSim & each actuator device & physics, viz, viewer & REST \\
IF-10 & Storage API & storage & controller, train, viz, people & REST \\
IF-11 & Room layers, alerts; graphs page & viz & viewer; a browser & REST \\
IF-12 & Device registry & registry & installer, devices, controller, viz & REST \\
\bottomrule
\end{longtable}

\subsection{The truth side}

\textbf{IF-1, clock and people.}
\code{GET occupancysim:8081/api/minutes?after=<time>} returns the clock and who
was in each room at every whole minute after \code{after}:

\begin{json}
{"clock": {"time": "2026-09-21T10:16:25Z", "factor": 120, "running": true},
 "minutes": [
  {"time": "2026-09-21T10:15:00Z", "rooms": {"level0/1570": {"occupants": 16}}},
  {"time": "2026-09-21T10:16:00Z", "rooms": {"level0/1570": {"occupants": 17, "guards": 1}}}]}
\end{json}

\code{occupants} counts everyone; \code{guards} says how many of them are night
guards. Empty rooms are left out. The last 240 minutes are kept. Physics asks
every 250~ms real. A paused clock gives no new minutes, so physics does not
step; it never extrapolates. On a timeout (500~ms) physics keeps its state and
asks again; the minutes it missed are still in the record. Physics also
publishes the time as a retained message on \code{time/model}
(\code{\{"model\_time": \ldots, "factor": 60\}}) every tick, which is where the
actuators take it from.

\textbf{IF-2, truth API.} \code{GET physics:8080/rooms?room=level0/A109}
returns \code{\{"run\_id", "model\_time", "rooms": \{"level0/A109":
\{"co2\_ppm": 612.4, "temp\_c": 21.3, "occupancy": 6\}\}\}}. Units: ppm, °C,
persons. \code{model\_time} is the whole minute the values belong to. A sensor
polls every 250~ms real and takes a reading when \code{model\_time} has
advanced by its sample period. On a timeout (500~ms) it publishes nothing, so
it goes silent downstream, which is what a dead device looks like.

\textbf{IF-5, truth topic} \code{truth/<level>/<room>}, one record per room and
model minute, readable by storage only:

\begin{json}
{"run_id": "eval-reactive-v2", "room": "level0/A109",
 "model_time": "2026-09-21T10:15:00Z",
 "co2_ppm": 612.4, "temp_c": 21.3, "occupancy": 6, "guards": 0,
 "airflow_ls": 180.0, "setpoint_c": 21.0,
 "fan_w": 270.0, "ahu_heat_w": 2160.0, "radiator_w": 310.0}
\end{json}

The three powers are averages over the minute that ended at \code{model\_time}
and are what the energy proxy sums. Physics does not buffer: a gap in the truth
shows as missing minutes in the evaluation and is reported, not filled.

\subsection{Observations}

\textbf{IF-3, readings in BuildSim.}
\code{PUT /api/sensors/\{device id\}-reading/value} with
\code{\{"data\_type": "text", "value": "612"\}}: \COtwo{} as an integer in ppm,
temperature with one decimal in °C, the count as an integer. One write per
sensor per sample (60~s). On failure the sensor retries three times 100~ms
apart, then drops the reading. A \code{404} means BuildSim has lost the
equipment; the registry recreates it within 10~s real.

\textbf{IF-4, observations} \code{obs/<level>/<room>/<kind>}, where
\code{kind} is \code{co2}, \code{temp} or \code{occupancy}, published right
after the IF-3 write:

\begin{json}
{"run_id": "eval-reactive-v2", "sensor_id": "co2-0017", "room": "level0/A109",
 "kind": "co2", "value": 612, "unit": "ppm",
 "model_time": "2026-09-21T10:15:00Z", "seq": 1789985700}
\end{json}

\code{seq} is the reading's model time in Unix seconds: unique per sensor,
increasing, and unchanged by a restart. Consumers drop duplicates on
\code{(sensor\_id, seq)}, because QoS~1 can deliver twice. A consumer judges
freshness by \code{model\_time}, never by arrival: older than 180~s is stale.
An injected fault is not marked in the message; finding it is the controller's
job. When no observation has arrived for 2.5 times the usual real-time gap (at
least 1.5~s), the controller and the dampers read IF-3 instead until messages
return.

\subsection{Decisions, commands and actuator state}

\textbf{IF-8, command.}
\code{POST http://<actuator endpoint>/rooms/\{level\}/\{room\}/command}. The
endpoint is the one the actuator advertised to the registry (IF-12).

\begin{json}
{"cmd_id": "1cb852f9fc5fe3f1579b2f0f5442e5b1", "value": 112.4, "unit": "l/s",
 "issued_at": "2026-09-21T09:28:00Z", "ttl_s": 300,
 "reason": "predictive: CO2 417 ppm, 11 people within 22 min"}
\end{json}

\begin{longtable}[]{@{}L{0.36\linewidth}L{0.56\linewidth}@{}}
\toprule
\textbf{Response} & \textbf{Meaning} \\
\midrule
\endhead
\code{202 \{"target": 112.4, "clamped": false\}} & Accepted. \code{target} is what the actuator will move to; \code{clamped} says the safety limits changed it \\
\code{200 \{"duplicate": true\}} & This \code{cmd\_id} was already handled; nothing changes \\
\code{409 \{"reason": "expired"\}} & \code{issued\_at + ttl\_s} is before the actuator's model time \\
\code{422 \{"reason": \ldots\}} & Out of range, wrong unit, unknown room, or no \code{cmd\_id} \\
\bottomrule
\end{longtable}

Units and ranges: a damper takes l/s from 0 to the room's maximum
(1.5 × design flow); a radiator valve takes a setpoint from 16 to 24~°C. The
controller sends at most one command per room and reading, waits 500~ms, and
retries the \emph{same} \code{cmd\_id} up to three times, which idempotency
makes safe. If no valid command arrives within \code{ttl\_s}, the actuator
returns to its fallback (design flow, or 21~°C) and publishes
\code{ttl\_expired}.

\textbf{IF-6, decision} \code{decision/<level>/<room>}, one record per command:

\begin{json}
{"run_id": "eval-predictive-v3", "room": "level0/1570",
 "model_time": "2026-09-21T09:28:00Z", "policy": "predictive",
 "mode": "normal", "fault": "", "observed": {"co2": 417},
 "predicted_occupancy": 11, "lead_min": 21.6, "airflow_target_ls": 112.4,
 "cmd_id": "1cb8...", "status": 202, "applied_target": 112.4, "clamped": false,
 "reason": "predictive: CO2 417 ppm, 11 people within 22 min"}
\end{json}

A heating decision carries \code{setpoint\_target\_c} instead of the airflow
fields. \code{mode} says what the decision rests on (\code{normal},
\code{no\_forecast}, \code{degraded\_occupancy}, \code{degraded\_co2},
\code{degraded\_temp}; Section 7.1), and \code{fault} says what the controller
saw when it found a sensor stuck, for example ``\COtwo{} sensor has read
425~ppm 5 times in a row''.

\textbf{IF-7, actuator events} \code{event/<level>/<room>}:
\code{\{"run\_id", "room", "model\_time", "actuator": "damper", "kind",
"detail"\}}, where \code{kind} is \code{override\_co2},
\code{override\_released}, \code{ttl\_expired} or \code{rejected}. Published
when the thing happens, not every tick.

\textbf{IF-9, actuator state in BuildSim.}
\code{PUT /api/actuators/\{device id\}-state/state} with
\code{\{"state": "180.0"\}}: l/s or °C with one decimal. The actuator writes
the state it has \emph{reached}, not the target. Airflow moves at most 10~\% of
the room's maximum per minute, a setpoint at most 0.5~°C per minute. The
actuator steps once per real second and writes when the value has changed.
Physics reads all states with one \code{GET /api/equipment?level=level0} per
tick; a missing or unreadable state means it keeps the last value.

\subsection{Storage, dashboard and registry}

\textbf{IF-10, storage API} (port 8090):

\begin{longtable}[]{@{}L{0.46\linewidth}L{0.46\linewidth}@{}}
\toprule
\textbf{Call} & \textbf{Returns} \\
\midrule
\endhead
\code{GET /history?run=\&room=\&kind=\&from=\&to=} & JSONL, one stored record per line, in the order stored. \code{kind} is \code{co2}, \code{temp}, \code{occupancy}, \code{truth}, \code{true\_occupancy}, \code{decision} or \code{event}; \code{from} is inclusive and \code{to} exclusive, in model time \\
\code{GET /runs} & The recorded run ids \\
\code{GET /runs/\{run\}/export} & A tar of the run's four JSONL files \\
\code{PUT /models/\{name\}}, \code{GET /models/latest} & Store a model; fetch the newest \\
\code{GET /healthz} & Counts of records stored, duplicates dropped and records rejected \\
\bottomrule
\end{longtable}

Storage checks every record before appending it: valid JSON, a safe
\code{run\_id}, a room and an RFC~3339 \code{model\_time}. A bad record is
counted and logged, never silently dropped. \code and
\code{kind=true\_occupancy} for the run in progress return \code{403}: the
oracle may only replay a finished run. Readers time out after 2~s; the
controller then keeps the model it has.

The model \code{train} writes is the per-room profile of Section 7.1:

\begin{json}
{"name": "profile-train-week-v3", "slot_min": 15, "quantile": 0.8,
 "slot_quantile": 0.75,
 "trained_on": {"run_ids": ["train-week-v2"], "weekdays": 5, "readings": 360000},
 "rooms": {"level0/1570": {"mean": [0, ...], "high": [0, ...], "days": [5, ...]}}}
\end{json}

\textbf{IF-11, dashboard.} viz is the only writer of room layers and alerts,
because every collection \code{PUT} in BuildSim replaces the whole collection.
\code{PUT /api/room-layers} sends all five layers once per real second;
\code{PUT /api/alerts} sends at most 100 alerts, most severe first, when the
list changes and every 30~s real. The graphs page is served on port 8060:
\code{GET /api/overview?minutes=N} for the building and
\code{GET /api/room?room=\ldots\&minutes=N} for one room, both with one value
per model minute and \code{null} where nothing was seen. Section 12 lists the
layers and alerts.

\textbf{IF-12, registry} (port 8070; the installer's form is at \code{/}):

\begin{longtable}[]{@{}L{0.40\linewidth}L{0.14\linewidth}L{0.38\linewidth}@{}}
\toprule
\textbf{Call} & \textbf{Who} & \textbf{Returns} \\
\midrule
\endhead
\code{POST /devices \{id, kind, room, installed\_by\}} & installer & \code{201} and the device; \code{409} id taken; \code{422} unknown kind, bad id, or a room that is not in the floor plan \\
\code{GET /devices?kind=\&room=\&status=} & controller, viz & The matching devices \\
\code{PUT /devices/\{id\}/checkin \{endpoint\}} & the device & \code{200} and its record; \code{404} not installed; \code{410} retired \\
\code{DELETE /devices/\{id\}} & installer & Retires the device and removes its equipment from BuildSim \\
\bottomrule
\end{longtable}

\begin{json}
{"id": "damper-0002", "kind": "damper", "room": "level0/1570", "status": "active",
 "installed_by": "kasper", "endpoint": "http://damper-0002:8080",
 "last_seen": "2026-09-30T18:21:12Z",
 "size": {"area_m2": 101.2, "volume_m3": 273.2, "capacity": 25,
          "design_ls": 210.4, "max_ls": 315.6,
          "occupied_min_ls": 35.4, "empty_min_ls": 10.1}}
\end{json}

A device checks in every 5~s real until it gets \code{200}; until then it
publishes nothing and accepts no commands. A device whose hardware kind differs
from the registered kind refuses to start. After that it checks in every 30~s
real, which is what \code{last\_seen} shows. The registry re-sends every
device's equipment to BuildSim every 10~s real, so a restarted BuildSim is
repaired without anyone noticing. Running devices keep working without the
registry; only starting devices wait for it.

\fussy
